% Fronts d'onde émis par une source en mouvement (L'essentiel, chapitre 23)
\begin{tikzpicture}[x=1cm,y=1cm]
\foreach \k in {1,...,5}{
  \pgfmathsetmacro\cx{-0.36*\k}\pgfmathsetmacro\r{0.6*\k}
  \draw[otNavy!70,line width=0.6pt] (\cx,0) circle (\r);}
\fill[otRed!80!black] (0,0) circle (0.09);
\draw[pcVitesse] (0.1,0) -- (0.8,0) node[pcLegende,anchor=south,text=otNavy] {$\vec{v}_\mathrm{E}$};
\node[pcLegende,anchor=north] at (0,-0.12) {S};
\fill[otInk] (2.0,0) circle (0.06) node[pcLegende,anchor=north] {A};
\fill[otInk] (-5.3,0) circle (0.06) node[pcLegende,anchor=north] {B};
\node[pcLegende,anchor=west,align=left] at (2.2,0.3) {fronts resserrés :\\ fréquence plus élevée};
\node[pcLegende,anchor=east,align=right] at (-5.5,0.3) {fronts espacés :\\ fréquence plus basse};
\end{tikzpicture}
\par\smallskip\legende{Chaque front d'onde est centré sur la position de la source à l'instant de son émission : la source rattrape les fronts qu'elle émet vers l'avant.}
